\documentclass[a4paper,12pt]{article}
\usepackage[ngerman]{babel}
\usepackage[utf8]{inputenc}
\usepackage[T1]{fontenc}
\usepackage{amsmath,amssymb}
\usepackage{tikz}
\usetikzlibrary{automata,positioning,arrows.meta}

\tikzset{
    ->,>=Stealth,
    node distance=2.5cm,
    every state/.style={thick,minimum size=1cm},
    initial text={},
}

\newcommand{\defbox}[1]{%
\medskip
\noindent\fbox{%
    \begin{minipage}{0.94\linewidth}
    \textbf{Definition / Vorgehen.}\par\smallskip
    #1
    \end{minipage}%
}
\medskip\par
}

\title{ETI -- Aufgabenblatt 2 \\ Hausaufgaben}
\author{}
\date{}

\begin{document}
\raggedright
\setlength{\parindent}{0pt}

{\LARGE\bfseries ETI -- Aufgabenblatt 2 \\ Hausaufgaben\par}
\vspace{2em}

\section*{Aufgabe 2.5}

\defbox{
Ein NFA akzeptiert ein Wort, wenn mindestens ein möglicher Berechnungspfad in einem akzeptierenden Zustand endet.
Die Sprache
\[
B=\{a,b\}^* \circ \{a\} \circ \{a,b\}^2
\]
besteht genau aus den Wörtern über $\{a,b\}$, deren drittes Zeichen von rechts ein $a$ ist. Für den NFA kann man dieses $a$ nichtdeterministisch raten. Für den DFA muss man deterministisch die letzten drei gelesenen Zeichen speichern.
}

\subsection*{1. NFA $M_1$ mit 4 Zuständen}

Der folgende NFA rät beim Lesen eines $a$, dass dieses $a$ das drittletzte Zeichen ist,
und liest danach noch genau zwei Zeichen.

\begin{flushleft}
\begin{tikzpicture}[node distance=2.8cm]
    \node[state, initial] (q0) {$q_0$};
    \node[state, right=of q0] (q1) {$q_1$};
    \node[state, right=of q1] (q2) {$q_2$};
    \node[state, accepting, right=of q2] (q3) {$q_3$};

    \path (q0) edge[loop above] node {$a,b$} ()
          (q0) edge node[above] {$a$} (q1)
          (q1) edge node[above] {$a,b$} (q2)
          (q2) edge node[above] {$a,b$} (q3);
\end{tikzpicture}
\end{flushleft}

\textbf{Korrektheit:}
Der Automat kann beliebig lange in $q_0$ bleiben. Sobald er ein $a$ liest, kann er nach $q_1$ wechseln und dieses $a$ als drittletztes Zeichen raten. Danach müssen noch genau zwei Zeichen gelesen werden, um in $q_3$ zu enden. Also akzeptiert $M_1$ genau die Wörter aus $\{a,b\}^*a\{a,b\}^2=B$.

\subsection*{2. DFA $M_2$ mit 8 Zuständen}

Der folgende DFA speichert immer die letzten drei Zeichen. Am Anfang werden fehlende Zeichen links mit $b$ aufgefüllt, damit Wörter der Länge $0,1,2$ nicht versehentlich akzeptiert werden.

\begin{center}
\begin{tikzpicture}[scale=0.82, transform shape]
    \node[state, accepting] (aaa) at (0,0) {$aaa$};
    \node[state, accepting] (aab) at (3,0) {$aab$};
    \node[state, accepting] (aba) at (6,0) {$aba$};
    \node[state, accepting] (abb) at (9,0) {$abb$};

    \node[state] (baa) at (0,-3) {$baa$};
    \node[state] (bab) at (3,-3) {$bab$};
    \node[state] (bba) at (6,-3) {$bba$};
    \node[state, initial] (bbb) at (9,-3) {$bbb$};

    \path (aaa) edge[loop above] node {$a$} ()
          (aaa) edge node[above] {$b$} (aab)
          (aab) edge node[above] {$a$} (aba)
          (aab) edge[bend left=18] node[above] {$b$} (abb)
          (aba) edge[bend right=18] node[below,sloped] {$a$} (baa)
          (aba) edge node[above,sloped] {$b$} (bab)
          (abb) edge node[above,sloped] {$a$} (bba)
          (abb) edge node[right] {$b$} (bbb)
          (baa) edge node[left] {$a$} (aaa)
          (baa) edge node[below,sloped] {$b$} (aab)
          (bab) edge node[below,sloped] {$a$} (aba)
          (bab) edge node[right] {$b$} (abb)
          (bba) edge[bend left=12] node[below] {$a$} (baa)
          (bba) edge[bend left=12] node[above] {$b$} (bab)
          (bbb) edge node[below] {$a$} (bba)
          (bbb) edge[loop below] node {$b$} ();
\end{tikzpicture}
\end{center}

\textbf{Korrektheit:}
Nach jedem gelesenen Wort $w$ enthält der Zustand die letzten drei Zeichen von $w$, wobei am Anfang links mit $b$ aufgefüllt wird. Ein Zustand ist genau dann akzeptierend, wenn sein erstes Zeichen $a$ ist. Das ist genau die Aussage, dass das drittletzte Zeichen des Eingabewortes ein $a$ ist. Also gilt $L(M_2)=B$.

\section*{Aufgabe 2.6}

\defbox{
Das Schubfachprinzip besagt: Werden $8$ Objekte auf höchstens $7$ Fächer verteilt, dann liegen zwei verschiedene Objekte im selben Fach. Für einen DFA bedeutet das: Wenn $8$ verschiedene Wörter nach höchstens $7$ Zuständen führen, dann müssen zwei verschiedene Wörter denselben Zustand erreichen. Wörter, die denselben Zustand erreichen, können durch dieselbe Fortsetzung nicht mehr unterschieden werden.
}

Wir beweisen per Widerspruch. Angenommen, es gibt einen DFA
\[
M=(Q,\{a,b\},\delta,q_0,F)
\]
mit $|Q|=7$ und $L(M)=B$.

Betrachte die $8$ Wörter aus $\{a,b\}^3$. Da $M$ nur $7$ Zustände hat, gibt es nach dem Schubfachprinzip zwei verschiedene Wörter
\[
x=x_1x_2x_3,\qquad y=y_1y_2y_3
\]
aus $\{a,b\}^3$, die nach dem Einlesen im selben Zustand landen:
\[
\delta^*(q_0,x)=\delta^*(q_0,y).
\]

Da $x\neq y$, unterscheiden sich die Wörter an mindestens einer Position.

\begin{itemize}
    \item Falls $x_1\neq y_1$, dann liegt genau eines von $x,y$ in $B$, denn bei Wörtern der Länge $3$ ist das erste Zeichen das drittletzte Zeichen.
    \item Falls $x_2\neq y_2$, dann liegt genau eines von $xa,ya$ in $B$, denn in $xa$ bzw. $ya$ ist das zweite Zeichen von $x$ bzw. $y$ das drittletzte Zeichen.
    \item Falls $x_3\neq y_3$, dann liegt genau eines von $xaa,yaa$ in $B$, denn in $xaa$ bzw. $yaa$ ist das dritte Zeichen von $x$ bzw. $y$ das drittletzte Zeichen.
\end{itemize}

In jedem Fall gibt es also eine Fortsetzung $z\in\{\varepsilon,a,aa\}$, so dass genau eines der Wörter $xz$ und $yz$ in $B$ liegt.
Weil $x$ und $y$ aber denselben Zustand erreichen, gilt
\[
\delta^*(q_0,xz)=\delta^*(q_0,yz).
\]
Der DFA müsste also beide Wörter gleich behandeln. Das widerspricht $L(M)=B$.

Damit kann es keinen DFA mit $7$ Zuständen geben, der $B$ erkennt.

\section*{Aufgabe 2.7}

\defbox{
Ein NFA akzeptiert ein Wort, wenn es mindestens einen Pfad vom Startzustand zu einem akzeptierenden Zustand gibt, dessen Kantenbeschriftungen genau das Wort ergeben. $\varepsilon$-Übergänge verbrauchen kein Eingabezeichen. Hier zerlegen wir den Automaten in drei Teilautomaten: den oberen, den mittleren und den unteren Ast.
}

Sei
\[
\Sigma=\{a,b,c\}.
\]
Der Startzustand ist $q_0$, die akzeptierenden Zustände sind $a_3,b_3,c_1$.

\begin{center}
\begin{tikzpicture}[node distance=2.2cm and 2.6cm]
    \node[state, initial] (q0) {$q_0$};

    \node[state, above=of q0] (a1) {$a_1$};
    \node[state, right=of a1] (a2) {$a_2$};
    \node[state, accepting, right=of a2] (a3) {$a_3$};

    \node[state, right=of q0] (b1) {$b_1$};
    \node[state, right=of b1] (b2) {$b_2$};
    \node[state, accepting, right=of b2] (b3) {$b_3$};

    \node[state, accepting, below=of q0] (c1) {$c_1$};

    \path (q0) edge[bend left=20] node[left] {$c$} (a1)
          (a1) edge node[above] {$a$} (a2)
          (a2) edge node[above] {$b$} (a3)
          (a1) edge[bend left=20] node[right] {$\varepsilon$} (q0)
          (a3) edge[out=-120,in=20] node[pos=0.63,above] {$\varepsilon$} (q0)
          (q0) edge node[above] {$\varepsilon$} (b1)
          (b1) edge[loop above] node {$a,b,c$} ()
          (b1) edge node[above] {$b$} (b2)
          (b2) edge[loop above] node {$b$} ()
          (b2) edge node[above] {$a,b,c$} (b3)
          (b3) edge[bend left=24] node[below] {$a,c$} (b1)
          (q0) edge[out=-110,in=110] node[pos=0.45,left] {$a$} (c1)
          (c1) edge[out=70,in=-70] node[pos=0.45,right] {$\varepsilon$} (q0);
\end{tikzpicture}
\end{center}

Wir betrachten den NFA als Zusammensetzung von drei Teilautomaten, die vom Auswahlzustand $q_0$ aus erreichbar sind.

\medskip
\textbf{Unterer Teilautomat:}
Der untere Ast liest genau
\[
C=a.
\]
Wenn der Automat in $c_1$ endet, wird akzeptiert. Über den $\varepsilon$-Übergang kann er aber auch wieder nach $q_0$ zurückkehren.

\medskip
\textbf{Oberer Teilautomat:}
Der obere Ast beginnt mit $c$. Danach kann er sofort per $\varepsilon$ nach $q_0$ zurückkehren oder mindestens einmal den Block $ab$ lesen und in $a_3$ akzeptieren. Damit gilt:
\[
A_0=c(ab)^*
\quad\text{als Rückkehr nach }q_0,
\]
\[
A_1=c(ab)^+
\quad\text{als akzeptierender Abschluss in }a_3.
\]

\medskip
\textbf{Mittlerer Teilautomat:}
Der mittlere Ast liest zuerst beliebig viele Zeichen, dann mindestens ein $b$ und danach noch ein beliebiges Zeichen aus $\Sigma$. Danach darf er mit $a$ oder $c$ wieder nach $b_1$ zurückspringen und denselben Teil erneut durchlaufen. Also ist
\[
B=\Sigma^*bb^*\Sigma\bigl((a\cup c)\Sigma^*bb^*\Sigma\bigr)^*.
\]

Vor dem letzten akzeptierenden Teil darf der Automat beliebig oft über den unteren oder oberen Ast nach $q_0$ zurückkehren. Die möglichen Rückkehr-Blöcke sind also
\[
R=C\cup A_0=a\cup c(ab)^*.
\]
Am Ende muss der Pfad in einem der drei akzeptierenden Teilautomaten enden. Daher ist die akzeptierte Sprache
\[
\boxed{
L=
R^*\bigl(C\cup A_1\cup B\bigr)
}
\]
also
\[
\boxed{
L=
\bigl(a\cup c(ab)^*\bigr)^*
\left(
a
\cup c(ab)^+
\cup \Sigma^*bb^*\Sigma\bigl((a\cup c)\Sigma^*bb^*\Sigma\bigr)^*
\right).
}
\]

\section*{Aufgabe 2.8}

\defbox{
Eine Sprache heißt regulär, wenn sie von einem DFA oder äquivalent von einem NFA erkannt wird. Für das Umkehren einer regulären Sprache verwendet man einen NFA: Man dreht alle Kanten um, macht die alten Endzustände zu möglichen Startzuständen und den alten Startzustand zum einzigen Endzustand.
}

Zu zeigen ist:
\[
L \text{ ist regulär}
\quad\Longleftrightarrow\quad
L^R \text{ ist regulär}.
\]

\subsection*{``$\Rightarrow$''}

Sei $L$ regulär. Dann gibt es einen NFA
\[
N=(Q,\Sigma,\delta,q_0,F)
\]
mit $L(N)=L$.

Wir konstruieren einen NFA
\[
N^R=(Q\cup\{q_s\},\Sigma,\delta_R,q_s,\{q_0\})
\]
für $L^R$. Dabei ist $q_s$ ein neuer Startzustand.

Die Übergänge von $N^R$ sind:
\begin{itemize}
    \item Für jeden alten Endzustand $f\in F$ gibt es einen $\varepsilon$-Übergang
    \[
    q_s \xrightarrow{\varepsilon} f.
    \]
    \item Für jeden Übergang $p\xrightarrow{x}q$ in $N$ mit $x\in\Sigma\cup\{\varepsilon\}$ gibt es in $N^R$ den umgedrehten Übergang
    \[
    q\xrightarrow{x}p.
    \]
\end{itemize}

Ein akzeptierender Pfad in $N$ mit Beschriftung
\[
w=w_1w_2\cdots w_n
\]
von $q_0$ zu einem Zustand $f\in F$ wird dadurch zu einem Pfad in $N^R$ von $q_s$ über $f$ nach $q_0$ mit Beschriftung
\[
w_n\cdots w_2w_1=w^R.
\]
Also erkennt $N^R$ genau $L^R$. Damit ist $L^R$ regulär.

\subsection*{``$\Leftarrow$''}

Sei nun $L^R$ regulär. Nach der gerade bewiesenen Richtung ist dann auch
\[
(L^R)^R
\]
regulär. Für jedes Wort gilt $(w^R)^R=w$, also ist
\[
(L^R)^R=L.
\]
Damit ist auch $L$ regulär.

\medskip
Folglich gilt: $L$ ist genau dann regulär, wenn $L^R$ regulär ist.

\end{document}
